% CAPÍTULO 1 INTRODUCCIÓN
% Protocolo: 2026-A135 — TT-2 Sistema Inteligente para la Identificación y Seguimiento de NNA en Situación de Orfandad por Feminicidio


\chapter*{INTRODUCCIÓN}
\addcontentsline{toc}{chapter}{INTRODUCCIÓN}

El presente Trabajo Terminal registrado bajo el número de protocolo 2026-A135, constituye la segunda fase de un proyecto de investigación aplicada que aborda una problemática
crítica en la intersección de los derechos humanos, la violencia de género y la ingeniería de sistemas inteligentes: la invisibilidad estadística y la ausencia de seguimiento
sistemático de las Niñas, Niños y Adolescentes (NNA) en situación de orfandad por feminicidio en México. Mediante la conceptualización, diseño, desarrollo y despliegue de un
sistema inteligente de monitoreo de medios, el proyecto propone una solución tecnológica para mitigar la dispersión de datos que obstaculiza la labor de las organizaciones
civiles y de las instituciones gubernamentales en la procuración de justicia y la reparación integral del daño.


\section*{Presentación del proyecto}
\label{sec:presentacion}


México atraviesa una crisis sistémica de violencia feminicida de proporciones documentadas. De acuerdo con los registros del Secretariado Ejecutivo del Sistema Nacional de Seguridad
Pública (SESNSP) y la Secretaría de Gobernación, entre los años 2010 y 2023 se han contabilizado \textbf{más de 8{,}000 víctimas} de feminicidio en el territorio nacional. Esta cifra
ya de por sí alarmante, oculta una dimensión adicional del daño: el impacto colateral sobre los hijos e hijas de las víctimas, quienes quedan en situación de orfandad repentina y traumática.

La Red por los Derechos de la Infancia en México (REDIM) estima que existen \textbf{al menos 3{,}500 Niñas, Niños y Adolescentes} que han perdido a su madre como consecuencia directa de un
feminicidio. Pese a la gravedad de esta cifra no existe a nivel nacional un padrón oficial, un censo unificado ni un mecanismo de seguimiento estandarizado que documente el estado de estas
víctimas indirectas. La Convención sobre los Derechos del Niño y la Ley General de los Derechos de Niñas, Niños y Adolescentes (LGDNNA) establecen el principio del \textit{interés superior
de la niñez} como eje que orienta, garantizando el derecho a la protección, la salud, la educación y la reparación integral. No obstante la ausencia de datos estructurados imposibilita la
materialización efectiva de estos derechos.

El TT no. 2026-A135 surge como respuesta tecnológica a esta brecha. El sistema automatiza el ciclo completo de descubrimiento informativo desde la recolección de noticias en fuentes periodísticas
de acceso público hasta la clasificación semántica y la persistencia relacional de los casos identificados, esto con el objetivo de que ningún NNA en situación de orfandad por feminicidio permanezca
invisible para las instituciones que tienen la obligación legal de protegerlo.


\section*{Planteamiento del problema}
\label{sec:planteamiento_problema}

La información sobre los NNA en situación de orfandad por feminicidio se encuentra fragmentada en el ecosistema informativo mexicano. Los datos relevantes sobre la existencia, el estado legal y las condiciones
de resguardo de los hijos de una víctima suelen estar dispersos en notas periodísticas de medios locales, comunicados aislados de fiscalías estatales y registros internos de organizaciones no gubernamentales
que no se interconectan entre sí. La información existe pero no es accesible de forma sistemática ni procesable a escala computacional.

\subsection*{La inviabilidad del monitoreo manual}
\label{subsec:monitoreo_manual}

La mitigación de esta invisibilidad estadística recae actualmente sobre organizaciones de la sociedad civil como la Fundación "Futuro con Derechos" el proyecto cartográfico de la investigadora
María Salguero y la REDIM. La metodología subyacente a este trabajo de documentación presenta un cuello de botella operativo: la \textbf{revisión manual y artesanal de fuentes hemerográficas}.
El proceso exige que analistas y voluntarios dediquen muchas horas semanales a la revisión individual de portales de noticias, buscando menciones colaterales de hijos en narrativas de notas policiacas.
Este enfoque es inviable frente al volumen de información que se genera diariamente en el país ya que México cuenta con más de 1{,}200 medios periodísticos digitales activos, produciendo decenas de
miles de artículos diarios.

La latencia introducida por el factor humano limita drásticamente la capacidad de procesamiento, provocando que una enorme cantidad de casos pasen desapercibidos porque es imposible auditar
concurrentemente todas las fuentes periodísticas nacionales y regionales. Además este proceso consume los recursos de talento especializado (psicólogos, abogados, trabajadores sociales),
desviando su capacidad hacia labores de recolección de datos en lugar de la intervención directa con las víctimas.


\section*{Objetivos del Proyecto}
\label{sec:objetivos}


\subsection*{Objetivo General}
\label{subsec:objetivo_general}

Desarrollar un sistema inteligente que sea capaz de automatizar la recolección, clasificación y análisis de 
información relacionadas con casos de feminicidio, para facilitar la identificación de niños, niñas y adolescentes 
que se encuentren en una situación de orfandad en consecuencia de los crímenes de violencias. Mediante el uso 
de Web Scraping y Procesamiento de Lenguaje Natural

\subsection*{Objetivos Específicos}
\label{subsec:objetivos_especificos}

\begin{description}
    \item[\textbf{OE-1}]
          Identificar y seleccionar fuentes públicas confiables que contengan información relacionada con feminicidios y posibles víctimas indirectas.

    \item[\textbf{OE-2}]
          Diseñar y desarrollar una herramienta automatizada de Web Scraping para la recolección sistemática de datos provenientes de dichas fuentes.

    \item[\textbf{OE-3} ]
          Implementar técnicas de Procesamiento de Lenguaje Natural (PLN) para extraer y estructurar información clave a partir de textos no estructurados.

    \item[\textbf{OE-4}]
          Diseñar una base de datos estructurada que permita almacenar la información recolectada de manera eficiente y segura.

\end{description}

